\begin{frame}{Overview}
\protect\phantomsection\label{overview}
\begin{block}{Where we are}
\protect\phantomsection\label{where-we-are}
Since the last meeting (16/04), the two main bugs were fixed and full
Combine results were obtained for the \(\mu\tau_h\) channel:

{\def\LTcaptype{none} % do not increment counter
\begin{longtable}[]{@{}ll@{}}
\toprule\noalign{}
Quantity & Value \\
\midrule\noalign{}
\endhead
Expected limit (median) & \(r < 2.04\) \\
Observed limit & \(r < 6.03\) \\
Observed significance & \(2.41\sigma\) \\
Expected significance & \(1.62\sigma\) \\
Best-fit \(r\) & \(2.13\;^{+1.92}_{-1.17}\) \\
GoF \(p\)-value & 0.51 \\
\bottomrule\noalign{}
\end{longtable}
}

\(\sigma_\text{SM} = 4.7\) fb, BDT Scan A (AdaBoost, \(N=600\),
\(d=2\)), per-arm proton selection.
\end{block}

\begin{block}{Today}
\protect\phantomsection\label{today}
\begin{enumerate}
\tightlist
\item
  Where the analysis stands and what is left
\item
  Translating \(r\) limits to \(a_\tau\) --- context with SM and
  existing constraints
\item
  Thesis outline and structure discussion
\end{enumerate}
\end{block}
\end{frame}

\begin{frame}[fragile]{Analysis status: what is done}
\protect\phantomsection\label{analysis-status-what-is-done}
\begin{block}{\(\mu\tau_h\) channel --- completed}
\protect\phantomsection\label{mutau_h-channel-completed}
\begin{itemize}
\tightlist
\item
  Phase 0 + Phase 1 pipeline (parallelised: days \(\to\) hours)
\item
  Normalization fixed: DY scale \(= 1.004\times10^{-4}\), PU proton
  weight \(= 0.245\)
\item
  Pileup proton treatment corrected (QCD stays data-driven, only DY and
  \(t\bar{t}\) enriched)
\item
  \textbf{Bug 1 fixed:} proton selection changed to per-arm logic (+25\%
  signal yield: \(14\,085 \to 21\,398\) events)
\item
  \textbf{Bug 2 fixed:} signal normalization corrected using
  \(\sigma_\text{SM} = 4.7\) fb and \(\sum w_\text{SM} = 402.661\) (was
  \(47\times\) too small)
\item
  BDT trained and scanned (Scan A selected); statistical validation
  complete
\end{itemize}
\end{block}

\begin{block}{Dominant uncertainty}
\protect\phantomsection\label{dominant-uncertainty}
\texttt{xiProton} (50\%) is the leading nuisance --- impacts \(r\) by
\(+49\,/\,-33\).
\end{block}
\end{frame}

\begin{frame}[fragile]{Analysis status: what is left}
\protect\phantomsection\label{analysis-status-what-is-left}
\begin{block}{Remaining items for \(\mu\tau_h\)}
\protect\phantomsection\label{remaining-items-for-mutau_h}
\begin{enumerate}
\tightlist
\item
  \textbf{Post-fit shape validation} --- identify which BDT bins drive
  the \(2.41\sigma\) excess; check pre/post-fit agreement per bin
\item
  \textbf{\(\sigma_\text{SM}\) cross-check} --- verify the signal MC
  normalization against an independent sample count
\item
  \textbf{Proton modelling} --- \texttt{xiProton} dominates; assess
  whether 50\% is the appropriate uncertainty or if it can be tightened
  with a dedicated study
\item
  \textbf{Background control regions} --- DY, \(t\bar{t}\), QCD
  normalization factors (\texttt{beta\_dy}, \texttt{beta\_tt},
  \texttt{beta\_qcd}) driven largely by the fit; cross-check with
  dedicated CRs
\end{enumerate}
\end{block}

\begin{block}{Extension}
\protect\phantomsection\label{extension}
\begin{itemize}
\tightlist
\item
  Apply the same pipeline to the \(\tau_h\tau_h\) channel (code
  structure already in place)
\item
  Combine \(\mu\tau_h\) + \(\tau_h\tau_h\) for improved sensitivity
\end{itemize}
\end{block}
\end{frame}

\begin{frame}[fragile]{From \(r\) limits to \(a_\tau\) --- the
connection}
\protect\phantomsection\label{from-r-limits-to-a_tau-the-connection}
\begin{block}{Signal cross section depends on \(a_\tau\)}
\protect\phantomsection\label{signal-cross-section-depends-on-a_tau}
The exclusive process \(\gamma\gamma \to \tau^+\tau^-\) probes the
\(\tau\tau\gamma\) vertex. The cross section receives a BSM correction:

\[\sigma(a_\tau, d_\tau) = \sigma_\text{SM} + \Delta\sigma(a_\tau - a_\tau^\text{SM}, d_\tau)\]

An observed (or excluded) signal strength
\(r = \sigma/\sigma_\text{SM}\) therefore constrains deviations
\(\Delta a_\tau = a_\tau - a_\tau^\text{SM}\).
\end{block}

\begin{block}{Strategy}
\protect\phantomsection\label{strategy}
\begin{itemize}
\tightlist
\item
  With Combine limits on \(r\), extract 95\% CL allowed range for
  \(a_\tau\) using the BSM reweighting branches already stored in the
  signal MC (\texttt{BSM\textbackslash{}\_weight} array)
\item
  Each entry in the array corresponds to a different
  \((a_\tau, d_\tau)\) hypothesis
\item
  Scan \(r(a_\tau)\) and read off where \(r_\text{obs/exp}\) intersects
  the curve
\end{itemize}
\end{block}
\end{frame}

\begin{frame}{\(a_\tau\): SM prediction and existing limits}
\protect\phantomsection\label{a_tau-sm-prediction-and-existing-limits}
\begin{block}{SM prediction}
\protect\phantomsection\label{sm-prediction}
\[a_\tau^\text{SM} = a_\tau^\text{QED} + a_\tau^\text{EW} + a_\tau^\text{HVP} + a_\tau^\text{HLbL} \approx 117\,721 \times 10^{-8}\]
\end{block}

\begin{block}{Existing experimental constraints (95\% CL)}
\protect\phantomsection\label{existing-experimental-constraints-95-cl}
{\def\LTcaptype{none} % do not increment counter
\begin{longtable}[]{@{}
  >{\raggedright\arraybackslash}p{(\linewidth - 4\tabcolsep) * \real{0.3750}}
  >{\raggedright\arraybackslash}p{(\linewidth - 4\tabcolsep) * \real{0.3438}}
  >{\raggedright\arraybackslash}p{(\linewidth - 4\tabcolsep) * \real{0.2812}}@{}}
\toprule\noalign{}
\begin{minipage}[b]{\linewidth}\raggedright
Experiment
\end{minipage} & \begin{minipage}[b]{\linewidth}\raggedright
Constraint
\end{minipage} & \begin{minipage}[b]{\linewidth}\raggedright
Process
\end{minipage} \\
\midrule\noalign{}
\endhead
L3 / OPAL (LEP) & \(-0.052 < a_\tau < 0.013\) &
\(Z \to \tau\tau\gamma\) \\
DELPHI (LEP) & similar order & \(e^+e^- \to e^+e^-\tau^+\tau^-\) \\
ATLAS (PbPb, 2023) & \(-0.0042 < a_\tau < 0.0048\) &
\(\gamma\gamma \to \tau\tau\) (UPC) \\
CMS (PbPb, 2023) & similar & \(\gamma\gamma \to \tau\tau\) (UPC) \\
\textbf{This analysis (pp)} & \textbf{TBD} &
\(\gamma\gamma \to \tau\tau\) + PPS tag \\
\bottomrule\noalign{}
\end{longtable}
}

The LEP limits are \(\sim 1\) order of magnitude above the first QED
contribution.\\
The PbPb UPC results are the current best, but pp collisions reach
\textbf{higher \(M_{\tau\tau}\)}, enhancing BSM sensitivity at large
dipole values.
\end{block}
\end{frame}

\begin{frame}{\(a_\tau\) landscape --- where this analysis sits}
\protect\phantomsection\label{a_tau-landscape-where-this-analysis-sits}
\begin{columns}
\column{0.55\textwidth}

### Why pp + PPS is complementary to PbPb

\begin{itemize}
\item PbPb UPC: lower $M_{\tau\tau}$ threshold, clean topology, large cross section
\item pp: higher $M_{\tau\tau}$ reach $\to$ BSM dipole effects scale as $M^2 / \Lambda^2$, so sensitivity to heavy new physics is enhanced
\item Proton tagging provides a strong kinematic constraint (mass matching) that is absent in PbPb
\item Expected pp limits on $a_\tau$ competitive with or beyond current LEP bounds
\end{itemize}

\column{0.45\textwidth}

### Sensitivity scaling

$$\Delta a_\tau^\text{BSM} \propto \frac{m_\tau^2}{\Lambda^2}$$

BSM sensitivity of $a_\tau$ is $\sim 280\times$ larger than $a_\mu$ and $\sim 3600\times$ larger than $a_e$ for the same NP scale $\Lambda$.

\medskip
$a_\tau$ is essentially **unmeasured** experimentally — the SM prediction has never been tested to the level of the first QED contribution.

\end{columns}
\end{frame}

\begin{frame}{Next step: \(r \to a_\tau\) conversion}
\protect\phantomsection\label{next-step-r-to-a_tau-conversion}
\begin{block}{Approach using existing BSM weights}
\protect\phantomsection\label{approach-using-existing-bsm-weights}
The signal MC already contains per-event BSM reweighting arrays. The
procedure:

\begin{enumerate}
\tightlist
\item
  For each hypothesis \(a_\tau^{(k)}\) in the scan, reweight the signal
  histogram:
  \[w_i^{(k)} = w_i^\text{SM} \times \frac{\text{BSM\_weight}[k]}{\text{BSM\_weight}[0]}\]
\item
  Rerun Combine (or use analytic scaling) to obtain
  \(r_{95}(a_\tau^{(k)})\)
\item
  Plot \(r_{95}\) vs \(a_\tau\) --- the crossing with
  \(r_\text{obs} = 6.03\) and \(r_\text{exp} = 2.04\) gives the excluded
  range
\end{enumerate}
\end{block}

\begin{block}{Goal for thesis}
\protect\phantomsection\label{goal-for-thesis}
Produce a plot showing: - Observed and expected excluded regions in
\(a_\tau\) - SM prediction marked - Comparison band with current best
limits (ATLAS/CMS PbPb)
\end{block}
\end{frame}

\begin{frame}{Thesis outline}
\protect\phantomsection\label{thesis-outline}
\begin{block}{Current chapter structure}
\protect\phantomsection\label{current-chapter-structure}
{\def\LTcaptype{none} % do not increment counter
\begin{longtable}[]{@{}
  >{\raggedright\arraybackslash}p{(\linewidth - 4\tabcolsep) * \real{0.2105}}
  >{\raggedright\arraybackslash}p{(\linewidth - 4\tabcolsep) * \real{0.3684}}
  >{\raggedright\arraybackslash}p{(\linewidth - 4\tabcolsep) * \real{0.4211}}@{}}
\toprule\noalign{}
\begin{minipage}[b]{\linewidth}\raggedright
\#
\end{minipage} & \begin{minipage}[b]{\linewidth}\raggedright
Title
\end{minipage} & \begin{minipage}[b]{\linewidth}\raggedright
Status
\end{minipage} \\
\midrule\noalign{}
\endhead
1 & Theory (literature review of key papers) & draft \\
2 & Introduction (SM, \(\tau\) lepton, \(g-2\), CEP motivation) & in
progress \\
3 & Experimental apparatus (CMS + PPS) & draft \\
4 & Preliminary analysis (conceptual stage, small sample validation) &
draft \\
5 & Samples and object reconstruction & in progress \\
6 & Analysis and results (\(\mu\tau_h\), BDT, Combine) & in progress \\
\bottomrule\noalign{}
\end{longtable}
}
\end{block}
\end{frame}

\begin{frame}{Thesis structure --- open questions}
\protect\phantomsection\label{thesis-structure-open-questions}
\begin{block}{Question 1: Chapter 1 placement}
\protect\phantomsection\label{question-1-chapter-1-placement}
Chapter 1 is currently a set of literature review notes on key papers.\\
\textbf{Option A:} Keep as standalone chapter (current) --- gives the
reader the bibliography context upfront.\\
\textbf{Option B:} Merge content into Chapter 2 (Introduction) ---
standard thesis structure; avoids a chapter that reads as notes.
\end{block}

\begin{block}{Question 2: Scope of results}
\protect\phantomsection\label{question-2-scope-of-results}
Currently only \(\mu\tau_h\) has full results.\\
\textbf{Option A:} Present \(\mu\tau_h\) as the main result with an
outlook section for other channels.\\
\textbf{Option B:} Attempt \(\tau_h\tau_h\) before submission and
include as a second result.
\end{block}

\begin{block}{Question 3: \(a_\tau\) limit extraction}
\protect\phantomsection\label{question-3-a_tau-limit-extraction}
Should the thesis include a full \(r \to a_\tau\) scan (using BSM
weights), or report limits on \(r\) only and leave \(a_\tau\)
translation as future work?
\end{block}
\end{frame}

\begin{frame}{Summary}
\protect\phantomsection\label{summary}
\begin{block}{Done since last meeting}
\protect\phantomsection\label{done-since-last-meeting}
Nothing new since 16/04 --- the corrected combine results stand as the
current best output.
\end{block}

\begin{block}{Immediate next steps}
\protect\phantomsection\label{immediate-next-steps}
\begin{enumerate}
\tightlist
\item
  Post-fit shape validation --- identify \(2.41\sigma\) excess origin
\item
  \(r \to a_\tau\) scan using BSM weight branches
\item
  Produce \(a_\tau\) exclusion plot vs SM + PbPb limits
\item
  Resolve thesis structure questions above
\item
  Extend pipeline to \(\tau_h\tau_h\)
\end{enumerate}
\end{block}

\begin{block}{Key numbers to keep in mind}
\protect\phantomsection\label{key-numbers-to-keep-in-mind}
\[r_\text{exp} < 2.04 \quad r_\text{obs} < 6.03 \quad \sigma_\text{SM} = 4.7\ \text{fb} \quad \text{GoF } p = 0.51\]
\end{block}
\end{frame}
